\documentclass[11pt]{article}

\usepackage[margin=1in]{geometry}
\usepackage{amsmath,amssymb,amsthm}
\usepackage{graphicx}
\usepackage{hyperref}

\newtheorem{lemma}{Lemma}
\newtheorem{theorem}{Theorem}
\newtheorem{corollary}{Corollary}

\newcommand{\HH}{\mathcal H}
\newcommand{\K}{\mathcal K}
\newcommand{\I}{\mathcal I}
\newcommand{\sa}{\mathrm{sa}}
\newcommand{\dom}{\operatorname{dom}}
\newcommand{\Int}{\operatorname{int}}
\newcommand{\Tr}{\operatorname{Tr}}

\title{Countable Spectral Extensions: One Positive Case and Empty-Interior Obstructions}
\author{arxiv\_automatic\_math run \texttt{fa\_banach\_001}}
\date{June 20, 2026}

\begin{document}
\maketitle

\section*{Source Question}

This packet concerns Remark 4.21(iii) in Ryshard-Pavel Kostecki,
\emph{Vainberg--Bregman relative entropy and quasinonexpansive operators},
arXiv:2511.14873 \cite{Kostecki}.

\begin{center}
\includegraphics[width=\textwidth]{figures/open_problem_crop.png}
\end{center}

The source asks for examples of unitarily invariant Euler--Legendre
functions $f\circ\widetilde\lambda$ on suitable separable reflexive
spaces of compact operators, and asks in particular whether this works
for the countable extensions of the finite spectral functions listed in
Proposition 4.17.

\section*{Claimed Status}

Candidate partial result, likely valid pending human review. The packet
classifies the direct countable extensions of the finite examples in
Proposition 4.17 into the one full-domain positive case and the
domain-restricted cases excluded by empty interior. The broader problem
of finding genuinely new non-$\Psi_\varphi$ examples remains open.

\section*{Proof Intuition}

The finite list has two different behaviours. The full-domain power
example is just the usual Legendre function $(1/p)\|x\|_p^p$ on
$\ell_p$ and on the Schatten class $S_p$, so it survives in infinite
dimension. The entropy examples whose domains are the positive cone, the
positive definite cone, or the operator interval $0\leq x\leq I$ do not
survive the same way. In infinite-dimensional compact-operator ideals,
every positive compact operator has arbitrarily small spectral
directions. A tiny negative rank-one perturbation in such a direction is
arbitrarily small in any compact ideal norm, but it exits the positive
cone. Thus those natural domains have empty norm interior, so they cannot
meet the nonempty-interior hypotheses used in Proposition 4.20.

\section*{Setup}

Let $\HH$ be an infinite-dimensional separable Hilbert space. Let
$\I(\HH)$ be any Banach ideal of compact operators on $\HH$ whose norm is
unitarily invariant and finite on rank-one operators; this includes the
Schatten classes and the rearrangement-invariant compact ideals appearing
in Proposition 4.20 of the source paper. Let
\[
  \I(\HH)^\sa=\{x\in\I(\HH):x=x^*\}
\]
with the ideal norm.

For a Borel scalar function $\phi$ on a real interval, the naive
countable spectral extension of the finite-dimensional functional
$\sum_i \phi(\lambda_i)$ has natural domain determined by the spectral
constraint on the compact self-adjoint operator. For the domain-restricted
examples in Proposition 4.17 and Example 2.5 of the source, these domains
are contained in one of:
\[
  \I(\HH)^+,\qquad
  \I(\HH)^+_0,\qquad
  \{x\in\I(\HH)^\sa:0\le x\le I\}.
\]
Here $\I(\HH)^+$ denotes positive compact operators in the ideal, and
$\I(\HH)^+_0$ denotes the strictly positive part of that cone when such a
condition is imposed by the formula.

\begin{lemma}[The full-domain power extension works]
Let $1<p<\infty$ and define $f:\ell_p\to\mathbb R$ by
\[
  f(x)=\frac1p\sum_{n=1}^\infty |x_n|^p=\frac1p\|x\|_p^p.
\]
Then $f$ is Euler--Legendre on $\ell_p$. Moreover, for the Schatten class
$S_p(\HH)^\sa$ one has
\[
  (f\circ\widetilde\lambda)(a)=\frac1p\|a\|_{S_p}^p,
\]
and $f\circ\widetilde\lambda$ is a unitarily invariant Euler--Legendre
function on $S_p(\HH)^\sa$.
\end{lemma}

\begin{proof}
Since $1<p<\infty$, both $\ell_p$ and its dual $\ell_q$ are reflexive,
strictly convex, and smooth. The subdifferential of
$f=(1/p)\|\cdot\|_p^p$ is the usual duality map
\[
  \nabla f(x)=(|x_n|^{p-2}x_n)_n\in\ell_q,
  \qquad q=\frac{p}{p-1},
\]
with the usual interpretation at $x_n=0$. Its Fenchel dual is
$f^*(y)=(1/q)\|y\|_q^q$, whose gradient is the inverse duality map.
Thus $f$ and $f^*$ are Gateaux differentiable on their whole spaces and
strictly convex; equivalently, $f$ is Euler--Legendre.

For $a\in S_p(\HH)^\sa$, $\widetilde\lambda(a)$ is the signed decreasing
rearrangement of the eigenvalues of $a$, counted with multiplicity, and
therefore
\[
  \|\widetilde\lambda(a)\|_{\ell_p}=\|a\|_{S_p}.
\]
Hence $(f\circ\widetilde\lambda)(a)=(1/p)\|a\|_{S_p}^p$. The same
argument as above, now applied to the uniformly convex and smooth
Schatten class $S_p$, shows that this function is Euler--Legendre on
$S_p(\HH)^\sa$. Unitary invariance follows from invariance of the
singular/eigenvalue list under unitary conjugacy.
\end{proof}

\begin{lemma}
The positive cone $\I(\HH)^+$ has empty interior in $\I(\HH)^\sa$.
\end{lemma}

\begin{proof}
Fix $a\in\I(\HH)^+$ and $\varepsilon>0$. Since $a$ is compact and
$\HH$ is infinite dimensional, $0$ lies in the spectrum of $a$ as an
accumulation point or eigenvalue. Hence for every $\delta>0$ there is a
unit vector $\xi\in\HH$ with $\langle a\xi,\xi\rangle<\delta$.

Let $p_\xi$ be the rank-one projection onto $\mathbb C\xi$. Since the
ideal norm is finite on rank-one operators, choose $\delta>0$ so small
that
\[
  \delta\|p_\xi\|_{\I}<\varepsilon
  \quad\text{and}\quad
  \langle a\xi,\xi\rangle<\delta .
\]
Then $b=a-\delta p_\xi$ belongs to $\I(\HH)^\sa$ and
$\|b-a\|_\I<\varepsilon$, but
\[
  \langle b\xi,\xi\rangle
  =\langle a\xi,\xi\rangle-\delta<0.
\]
Thus $b$ is not positive. Every ball around $a$ therefore meets the
complement of $\I(\HH)^+$.
\end{proof}

\begin{corollary}
The sets $\I(\HH)^+_0$ and
$\{x\in\I(\HH)^\sa:0\le x\le I\}$ also have empty interior in
$\I(\HH)^\sa$.
\end{corollary}

\begin{proof}
Both sets are contained in $\I(\HH)^+$, whose interior is empty.
\end{proof}

\begin{theorem}
Let $\HH$ be infinite dimensional. For the direct countable extensions of
the finite examples listed in Proposition 4.17 of the source paper, the
following holds.
\begin{itemize}
\item The full-domain power example
  $\gamma\sum_i |x_i|^{1/\gamma}$, with $p=1/\gamma\in(1,\infty)$,
  gives the positive example $(1/p)\|a\|_{S_p}^p$ on $S_p(\HH)^\sa$.
\item The following domain-restricted examples cannot satisfy
  Proposition 4.20's nonempty-interior hypotheses in any
  infinite-dimensional compact-operator ideal:
\begin{itemize}
\item the Kullback--Leibler/Umegaki entropy
  $\Tr(x\log x-x)$ on the positive compact cone;
\item the Burg/Pinsker logarithmic barrier $-\Tr(\log x)$ on the
  strictly positive compact cone;
\item the Fermi--Dirac entropy
  $\Tr(x\log x+(I-x)\log(I-x))$ on $0\le x\le I$;
\item the positive-domain Tsallis/power examples $\Tr(x^\alpha)$ with
  the positive-cone domain restrictions appearing in the finite list.
\end{itemize}
\end{itemize}
\end{theorem}

\begin{proof}
The positive full-domain statement is the preceding lemma, after the
identification $p=1/\gamma$.

Each listed domain-restricted countable extension has domain contained in
one of the empty-interior sets from the lemma and corollary. Consequently
\[
  \Int\dom(f\circ\widetilde\lambda)=\varnothing
\]
in the norm topology of $\I(\HH)^\sa$. Proposition 4.20 assumes, among
other things, that the relevant interiors of
$f\circ\widetilde\lambda$ and its Fenchel dual are nonempty. Therefore
these direct countable extensions cannot satisfy the setup of
Proposition 4.20.
\end{proof}

\section*{Scope and Limitations}

This answers only the direct countable-extension reading of the finite
list in Proposition 4.17. It does not rule out:
\begin{itemize}
\item other full-domain countable spectral functions beyond the norm-power
  case isolated above;
\item different topologies or non-norm interiors;
\item renormalised domains that are not contained in the positive compact
  cone or compact operator interval;
\item other unitarily invariant Euler--Legendre functions with
  $\Psi\ne\Psi_\varphi$.
\end{itemize}

\section*{Novelty and Literature Check}

The cheap run indexes were searched for arXiv:2511.14873, Bregman,
Vainberg, relative entropy, Euler--Legendre, and spectral examples, with
no duplicate packet found. A web search for the paper title and phrases
around ``countable extensions'' and ``Euler--Legendre Schatten'' found
the source arXiv page but no later explicit answer to Remark 4.21(iii).

\section*{Human Review Recommendation}

Review whether the source author's phrase ``countable extensions'' could
mean a non-norm topology or a renormalised domain. Under the norm topology
and the natural compact-operator domains displayed in Proposition 4.17,
the empty-interior obstruction is decisive.

\begin{thebibliography}{9}

\bibitem{Kostecki}
R.-P. Kostecki,
\emph{Vainberg--Bregman relative entropy and quasinonexpansive operators},
arXiv:2511.14873, 2025.

\end{thebibliography}

\end{document}
